% Part: history
% Chapter: biographies 
% Section: alfred-tarski

\documentclass[../../../include/open-logic-section]{subfiles}

\begin{document}

\olfileid{his}{bio}{tar}

\olsection{ॲल्फ्रेड टार्स्की}

\olphoto{tarski-alfred}{ॲल्फ्रेड टार्स्की}

ॲल्फ्रेड टार्स्की यांचा जन्म १४ जानेवारी १९०१ रोजी पोलंडमधील वॉर्सा येथे झाला; तेव्हा ते शहर रशियन साम्राज्याचा भाग होते. टार्स्की यांचे वर्णन “नेपोलियनसारखे” असे केले जात असे: ते उत्साही, बोलके आणि तीव्र स्वभावाचे होते. त्यांची ऊर्जा व्याख्यानांतही दिसत असे. एकदा व्याख्यानादरम्यान सिगारेट टाकताना त्यांनी कचरापेटीला आग लावली आणि त्यानंतर त्या इमारतीत त्यांना पुन्हा व्याख्यान देण्यास बंदी घालण्यात आली.

लहानपणापासूनच टार्स्की यांना ज्ञानाची तीव्र ओढ होती. आयुष्याच्या पुढील काळात ते विद्यार्थ्यांना सांगत की तर्कशास्त्रात त्यांना 'B' गुण मिळाले म्हणून त्यांनी तो विषय शिकला; तोच त्यांचा 'B' गुण मिळालेला एकमेव विषय होता. मात्र त्यांच्या माध्यमिक शाळेच्या नोंदींनुसार तर्कशास्त्रासह सर्व विषयांत त्यांना 'A' गुण मिळाले होते. त्यांनी १९१८ ते १९२४ या काळात वॉर्सा विद्यापीठात शिक्षण घेतले. सुरुवातीला त्यांचा जीवशास्त्र शिकण्याचा विचार होता; पण विद्यापीठ हे वॉर्सा तर्कशास्त्र व तत्त्वज्ञान परंपरेचे केंद्र असल्याने त्यांना गणित, तत्त्वज्ञान आणि तर्कशास्त्रात रस निर्माण झाला. १९२४ मध्ये त्यांनी स्तानिस्वाव लेश्न्येव्स्की (Stanis\l{}aw Le\'{s}niewski) यांच्या मार्गदर्शनाखाली डॉक्टरेट मिळवली.

१९३९ मध्ये अमेरिकेत स्थलांतर करण्यापूर्वी टार्स्की यांनी वॉर्सा येथे माध्यमिक शाळेत शिकवत असताना आपले काही सर्वात महत्त्वाचे काम पूर्ण केले. तार्किक परिणाम आणि तार्किक सत्य यांवरील त्यांचे लेख याच काळात लिहिले गेले. १९३९ मध्ये ते व्याख्यानांच्या दौऱ्यासाठी अमेरिकेत गेले होते. त्याच भेटीदरम्यान जर्मनीने पोलंडवर आक्रमण केले; ज्यू वंशामुळे टार्स्की परत जाऊ शकले नाहीत. त्यांची पत्नी आणि मुले युद्ध संपेपर्यंत पोलंडमध्ये राहिली; नंतर त्यांनाही अमेरिकेत स्थलांतर करता आले. टार्स्की यांनी हार्वर्ड, सिटी कॉलेज ऑफ न्यू यॉर्क, प्रिन्स्टन येथील इन्स्टिट्यूट फॉर ॲडव्हान्स्ड स्टडी आणि अखेरीस कॅलिफोर्निया विद्यापीठ, बर्कली येथे अध्यापन केले. बर्कलीमध्ये त्यांनी तर्कशास्त्र आणि विज्ञानपद्धती यांचा बहुविषयक अभ्यास कार्यक्रम स्थापन केला. २६ ऑक्टोबर १९८३ रोजी वयाच्या ८२व्या वर्षी टार्स्की यांचे निधन झाले.

\begin{reading}
टार्स्की यांच्या जीवनाविषयी अधिक माहितीसाठी \emph{Alfred Tarski: Life
  and Logic} हे चरित्र \citep{Feferman2004} पाहा. तार्किक परिणाम आणि सत्य यांवरील टार्स्की यांचे मूलभूत लेख इंग्रजीत \citep{Tarski1983} येथे उपलब्ध आहेत. त्यांचे सर्व मूळ लेख चार खंडांच्या मालिकेत संकलित केले आहेत \citep{Tarski1981}.
\end{reading}

\end{document}
